\begin{filecontents*}[overwrite]{algebraic_monodromy_verify.py}
"""Exact finite checks; these do not replace the proofs in the manuscript."""
from pathlib import Path
import json
import sympy as sp


def mul(a, b):
    return tuple(a[b[i]] for i in range(len(a)))


def inv(a):
    result = [0] * len(a)
    for i, image in enumerate(a):
        result[image] = i
    return tuple(result)


def generated(generators):
    identity = tuple(range(len(generators[0])))
    elements = {identity}
    todo = [identity]
    while todo:
        a = todo.pop()
        for b in generators:
            c = mul(a, b)
            if c not in elements:
                elements.add(c)
                todo.append(c)
    return sorted(elements)


def subgroup(generator):
    identity = tuple(range(len(generator)))
    result = [identity]
    g = generator
    while g != identity:
        result.append(g)
        g = mul(g, generator)
    return result


def cosets(group, inertia):
    return sorted({tuple(sorted(mul(h, a) for a in inertia)) for h in group})


def fixed_cosets(a, quotient):
    return sum(tuple(sorted(mul(a, h) for h in c)) == c for c in quotient)


def cover_check(name, generators, characters):
    group = generated(generators)
    order = len(group)
    index = {a: i for i, a in enumerate(group)}
    identity = tuple(range(len(generators[0])))
    quotients = [cosets(group, subgroup(a)) for a in generators]
    r = len(generators)
    rows = order + sum(map(len, quotients))
    edges = r * order
    boundary = sp.zeros(rows, edges)
    offset = order
    for i, quotient in enumerate(quotients):
        coset_index = {h: j for j, c in enumerate(quotient) for h in c}
        for h in group:
            column = i * order + index[h]
            boundary[index[h], column] = -1
            boundary[offset + coset_index[h], column] = 1
        offset += len(quotient)
    kernel_vectors = boundary.nullspace()
    dimension = len(kernel_vectors)
    kernel = (sp.Matrix.hstack(*kernel_vectors)
              if kernel_vectors else sp.zeros(edges, 0))
    pivots = list(kernel.T.rref()[1]) if dimension else []
    pivot_inverse = kernel[pivots, :].inv() if dimension else sp.zeros(0, 0)
    product = identity
    for a in generators:
        product = mul(product, a)
    infinity = cosets(group, subgroup(inv(product)))
    affine_character = {}
    projective_character = {}
    for a in group:
        permutation = []
        a_inverse = inv(a)
        for i in range(r):
            permutation.extend(i * order + index[mul(a_inverse, h)] for h in group)
        moved = kernel[permutation, :]
        coordinates = pivot_inverse * moved[pivots, :] if dimension else sp.zeros(0, 0)
        assert moved == kernel * coordinates
        trace = sp.trace(coordinates)
        expected = ((r - 1) * (order if a == identity else 0)
                    - sum(fixed_cosets(a, q) for q in quotients) + 1)
        assert trace == expected
        affine_character[a] = trace
        projective_character[a] = trace - fixed_cosets(a, infinity) + 1
    multiplicities = {}
    for label, character in characters.items():
        aff = sp.simplify(sum(affine_character[a] * sp.conjugate(character(a))
                              for a in group) / order)
        proj = sp.simplify(sum(projective_character[a] * sp.conjugate(character(a))
                               for a in group) / order)
        assert aff.is_integer and aff >= 0
        assert proj.is_integer and proj >= 0
        multiplicities[label] = {"affine": int(aff), "projective": int(proj)}
    compact_dimension = projective_character[identity]
    assert compact_dimension.is_integer and compact_dimension >= 0
    assert compact_dimension % 2 == 0
    assert sum(value["affine"] * characters[label](identity)
               for label, value in multiplicities.items()) == dimension
    assert sum(value["projective"] * characters[label](identity)
               for label, value in multiplicities.items()) == compact_dimension
    return {"name": name, "group_order": order, "finite_branch_count": r,
            "affine_homology_dimension": dimension,
            "genus": int(compact_dimension / 2),
            "points_above_infinity": len(infinity),
            "multiplicities": multiplicities,
            "deck_actions_checked": order}


records = []
# Rational characters suffice for C2, C2 x C2, and S3; no floating arithmetic.
c2 = (1, 0)
char_c2 = {"trivial": lambda a: sp.Integer(1),
           "sign": lambda a: sp.Integer(1 if a == (0, 1) else -1)}
for count in (1, 2, 3, 4, 5):
    records.append(cover_check(f"C2_{count}_finite_branches", [c2] * count, char_c2))

g1 = tuple(i ^ 1 for i in range(4))
g2 = tuple(i ^ 2 for i in range(4))
chars_biquad = {}
for bit in range(4):
    chars_biquad[f"chi_{bit}"] = (
        lambda a, bit=bit: sp.Integer((-1) ** ((bit & a[0]).bit_count())))
records.append(cover_check("biquadratic", [g1, g2], chars_biquad))

s = (1, 0, 2)
t = (0, 2, 1)


def parity(a):
    return (-1) ** sum(a[i] > a[j] for i in range(len(a)) for j in range(i + 1, len(a)))


chars_s3 = {"trivial": lambda a: sp.Integer(1),
            "sign": lambda a: sp.Integer(parity(a)),
            "standard": lambda a: sp.Integer(sum(a[i] == i for i in range(3)) - 1)}
records.append(cover_check("S3_two_transpositions", [s, t], chars_s3))
records.append(cover_check("S3_three_transpositions", [s, t, s], chars_s3))

# Cyclic groups C_m: compare the general character multiplicity formula with
# exact ranks of projection onto inertia-fixed vectors, represented as integers.
cyclic = []
for m in range(2, 13):
    for proposed_exponents in ((1,), (1, 1), (1, 2)):
        # A residue-zero inertia exponent is unramified, so it is not a
        # branch value. Removing that unnecessary star arm leaves both
        # homology character formulas unchanged.
        branch_exponents = tuple(a % m for a in proposed_exponents if a % m)
        infinity_exponent = -sum(branch_exponents) % m
        entries = []
        for character in range(m):
            fixed = [int(character * a % m == 0) for a in branch_exponents]
            fixed_infinity = int(character * infinity_exponent % m == 0)
            affine = len(branch_exponents) - 1 - sum(fixed) + int(character == 0)
            projective = affine - fixed_infinity + int(character == 0)
            assert affine >= 0 and projective >= 0
            entries.append((affine, projective))
        assert sum(p for _, p in entries) % 2 == 0
        cyclic.append({"degree": m, "finite_inertia_exponents": list(branch_exponents),
                       "affine_multiplicities": [a for a, _ in entries],
                       "projective_multiplicities": [p for _, p in entries]})

# Exact rational certificates for the Laurent cover x=t^m+t^(-m).
z = sp.symbols("t")
rational_checks = 0
for m in range(2, 7):
    dx = m * z ** (m - 1) - m * z ** (-m - 1)
    for k in range(-7, 8):
        if k % m == 0:
            continue
        primitive = sp.Rational(m, k + m) * z ** (k + m) - sp.Rational(m, k - m) * z ** (k - m)
        assert sp.expand(sp.diff(primitive, z) - z ** k * dx) == 0
        rational_checks += 1

# Polynomial inverse tower, checked directly as rational functions for 20 steps.
current = z
dx_cubic = 3 * z ** 2 + 1
for depth in range(1, 21):
    following = sp.integrate(sp.expand(current * dx_cubic), z)
    assert sp.expand(sp.diff(following, z) - dx_cubic * current) == 0
    current = following
    rational_checks += 1

result = {"arithmetic": "exact integer/rational symbolic arithmetic",
          "graph_cases": records, "cyclic_cases": cyclic,
          "rational_primitive_certificates_checked": rational_checks,
          "failures": 0,
          "scope": "Finite checks of characters and example certificates, not a formal verification of the general theorem."}
target = Path(__file__).with_suffix(".json")
target.write_text(json.dumps(result, indent=2) + "\n")
print(json.dumps({"graph_cases": len(records), "cyclic_cases": len(cyclic),
                  "deck_actions_checked": sum(r["deck_actions_checked"] for r in records),
                  "rational_primitive_certificates_checked": rational_checks,
                  "failures": 0}, indent=2))
\end{filecontents*}
\begin{filecontents*}[overwrite]{integration_verify.py}
"""Exact checks for the curve criterion and one-primitive residue theorem.

Finite nondependence checks are decisive only because the proved 2d/d+1
degree theorem supplies a finite stopping bound. This script checks examples,
not the general algebraic-geometric proof.
"""
import json
from pathlib import Path
import sympy as s

x, y, z, t = s.symbols('x y z t')
results=[]

def eq(name, lhs, rhs=0):
    residual=s.cancel(lhs-rhs)
    assert residual==0, (name,residual)
    results.append({'name':name,'pass':True})

def rank_over_k(fs, variable):
    den=s.prod(s.denom(s.cancel(f)) for f in fs)
    ps=[s.Poly(s.cancel(f*den),variable) for f in fs]
    n=max(p.degree() for p in ps)
    return s.Matrix([[p.nth(i) for p in ps] for i in range(n+1)]).rank()

def orbit(name,r,v,var,bound,expected_rank):
    delta=lambda f:s.cancel(s.diff(f,x)+v*s.diff(f,var))
    fs=[r]
    for j in range(bound): fs.append(delta(fs[-1]))
    rank=rank_over_k(fs,var)
    assert rank==expected_rank,(name,rank)
    results.append({'name':name,'pass':True,'derivatives_tested':bound+1,
                    'rank_over_Qx':rank,'proved_order_bound':bound})
    return fs

# Rational function fields are regular over k=Q(x), so bound d+1 applies.
orbit('Laurent eigenfunctions: degree 3, exact order 2',y+y**-2,y,y,4,2)
orbit('Ramified ambient field: r=y^2+y^4, exact order 3',y**2+y**4,1/(2*y),y,5,3)
orbit('Ramification obstruction: r=y, delta y=1/(2y)',y,1/(2*y),y,2,3)
orbit('Coupled pole: alpha=x+y, a=1/y, residue y',y,y/(1-y),y,2,3)
orbit('Moving residue pole with delta y=0',1/(y-x),s.Integer(0),y,2,3)
orbit('Pole relocation: delta y=y^2, r=1/y, exact order 2',1/y,y**2,y,2,2)

# Parameter-dependent primitive fields from closed rational one-forms.
for aa,bb,alpha,label in [
    (1/(x+y),1/(x+y),x*y,'log(x+y), alpha=xy'),
    (2*y/(x*x+y*y),2*x/(x*x+y*y),x+y*y,'log(x^2+y^2), alpha=x+y^2')]:
    eq('Closed one-form: '+label,s.diff(aa,x),s.diff(bb,y))
    ee=s.cancel(aa-s.diff(alpha,y))
    nu=s.cancel(s.diff(alpha,x)-bb)
    vv=s.cancel(nu/ee)
    dd=lambda h:s.cancel(s.diff(h,x)+vv*s.diff(h,y))
    eq('Parameter-dependent normalization: '+label,dd(ee)/ee+s.diff(vv,y))
    cc=x+y*y
    eq('Parameter-dependent simple residue transport: '+label,
       (s.diff(cc,x)+s.diff(vv*cc,y))/ee,dd(cc/ee))

aa=1/(x+y)
kernel=1/t
dxlog=lambda h:s.diff(h,x)+aa*s.diff(h,t)
dylog=lambda h:s.diff(h,y)+aa*s.diff(h,t)
eq('log(x+y) kernel telescope',dxlog(kernel),dylog(kernel))
orbit('log(x^2+y^2) kernel nonexistence by degree-two bound',
      (x*x+y*y)/(2*y),-x/y,y,3,4)

# Coupled primitive pole with a rational part in the primitive derivative.
for alpha, label in [(x*y,'xy'),(x+y*y,'x+y^2')]:
    aa=1/y
    ee=s.cancel(aa-s.diff(alpha,y))
    vv=s.cancel(s.diff(alpha,x)/ee)
    dd=lambda h:s.cancel(s.diff(h,x)+vv*s.diff(h,y))
    eq('Normalization identity for alpha='+label,
       dd(ee)/ee+s.diff(vv,y))
    cc=x+y*y
    raw=s.diff(cc,x)+s.diff(vv*cc,y)
    eq('Simple residue transport for alpha='+label,raw/ee,dd(cc/ee))
    # A double-pole term reduces to S(c)/p; test parameter transport.
    SS=lambda h:s.cancel(s.diff(h/ee,y))
    residue=s.cancel(SS(cc)/ee)
    # Dx(c/p^2) = c_x/p^2 + 2 alpha_x c/p^3.
    next_residue=s.cancel((SS(s.diff(cc,x))+SS(SS(2*s.diff(alpha,x)*cc))/2)/ee)
    eq('Double residue transport for alpha='+label,next_residue,dd(residue))

a=1+1/y
dy=lambda f:s.cancel(s.diff(f,y)+a*s.diff(f,t))
dx=lambda f:s.diff(f,x)
p=t-x-y
f=(1+y**-3)/p  # eta*(y+y^-2)/p, eta=1/y
P=lambda f:s.cancel(dx(dx(f))+dx(f)-2*f)  # (Dx-1)(Dx+2)
g=-((dx(y/p)+2*y/p)+(dx(y**-2/p)-y**-2/p))
eq('Coupled positive example: exact order-two certificate',P(f),dy(g))

# A genuinely algebraic coupled pole alpha^2=y+x. Work in k(z), y=z^2-x.
yp=z*z-x
a=1/yp
dyz=lambda f:s.cancel(s.diff(f,z)/(2*z)+a*s.diff(f,t))
dxz=lambda f:s.cancel(s.diff(f,x)+s.diff(f,z)/(2*z))
eta=a-1/(2*z)
f=(-t+2*z*z/yp)/(t*t-z*z)
c1=s.cancel(s.limit((t-z)*f,t,z))
eq('Algebraic pole normalized residue equals z',c1/eta,z)
df=dxz(f)
regular=s.cancel((t-z)**2*df)
c2=s.cancel(regular.subs(t,z))
c1new=s.cancel(s.diff(regular,t).subs(t,z))
transport=s.cancel((c1new+s.diff(c2/eta,z)/(2*z))/eta)
delta_z=1/(2*z-z*z+x)
eq('Algebraic pole residue transport',transport,delta_z)
vv=s.cancel((1/(2*z))/eta)
delta_coeff=lambda h:s.cancel(dxz(h)+vv*dyz(h))
eq('Algebraic pole normalization identity via implicit derivatives',
   delta_coeff(eta)/eta+s.diff(vv,z)/(2*z))
orbit('Algebraic coupled pole obstruction, normalized residue z',z,delta_z,z,2,3)

q=t*t-y-x
alog=1/y
dylog=lambda f:s.cancel(s.diff(f,y)+alog*s.diff(f,t))
f_log=dylog(q)/q
eq('Algebraic coupled positive example: logarithmic derivative',s.diff(f_log,x),dylog(-1/q))

# Polynomial remainder is treated by three rational-CT induction steps.
dy0=lambda f:s.cancel(s.diff(f,y)+s.diff(f,t)/y)
P1=lambda f:s.cancel(s.diff(f,x)+f/x)
P2=lambda f:s.cancel(s.diff(f,x)+2*f/x)
H=t*t/(x+y)
g2=1/(x+y)
R1=s.cancel(s.diff(H,x)-dy0(g2*t*t))
eq('Polynomial induction first remainder',R1,-2*t/(y*(x+y)))
g1=2/(x*(x+y))
R0=s.cancel(P1(R1)-dy0(g1*t))
eq('Polynomial induction second remainder',R0,-2/(x*y*(x+y)))
g0=2/(x*x*(x+y))
eq('Polynomial induction final rational identity',P2(R0),dy0(g0))
cert=s.cancel(P2(P1(g2*t*t))+P2(g1*t)+g0)
eq('Polynomial induction combined order-three telescope',P2(P1(s.diff(H,x))),dy0(cert))

out=Path(__file__).with_name('integration_verification.json')
out.write_text(json.dumps({'all_passed':True,'checks':results},indent=2))
print(json.dumps({'all_passed':True,'number_of_checks':len(results),'result_file':str(out)}))
\end{filecontents*}
\begin{filecontents*}[overwrite]{beta_problem3_certify.py}
"""Rigorous compact positivity certificate for Qi--Lim--Nantomah Problem 3.

All intervals have rational endpoints.  A comparison succeeds only when Arb
certifies the sign for the entire interval; there are no floating grid tests.
"""
from flint import arb, fmpq, ctx
from math import factorial

ctx.prec = 192

def iv(a, b):
    return arb(a).union(arb(b))

def pg(n, x):
    if n == 0:
        return x.digamma()
    return (-1)**(n+1) * factorial(n) * arb(n+1).zeta(x)

def data(x):
    z, u, v, w = [pg(k, x) for k in range(4)]
    Z, U, V = [pg(k, 1/x) for k in range(3)]
    H = z+x*u
    Hp = 2*u+x*v
    Hpp = 3*v+x*w
    return z,u,v,w,Z,U,V,H,Hp,Hpp

def cleared_log_convexity(x):
    z,u,v,w,Z,U,V,H,Hp,Hpp = data(x)
    return (x*x*((x*Hp+x*x*Hpp)*H-x*x*Hp*Hp)*z*z*Z*Z
            -x*x*((x*u+x*x*v)*z-x*x*u*u)*H*H*Z*Z
            -((x*U+V)*Z-U*U)*H*H*z*z)

def B_and_critical(x):
    z,u,v,w,Z,U,V,H,Hp,Hpp = data(x)
    P=z*Z
    Pp=u*Z-z*U/(x*x)
    ratio=(P/H).log()
    B=ratio/x.log()
    # sign of B', denominator x*(log x)^2 positive
    C=x*x.log()*(Pp/P-Hp/H)-ratio
    return B,C

def prove_positive(a,b):
    stack=[(fmpq(a),fmpq(b),0)]
    count=0
    certified=[]
    while stack:
        lo,hi,depth=stack.pop()
        val=cleared_log_convexity(iv(lo,hi))
        if val>0:
            count+=1
            certified.append((str(lo),str(hi),str(val)))
            continue
        if depth>=30:
            raise ArithmeticError((lo,hi,val))
        mid=(lo+hi)/2
        stack.extend([(lo,mid,depth+1),(mid,hi,depth+1)])
    print('CERTIFIED cleared log-convexity positive on [2/3,1]; pieces =',count)
    from pathlib import Path
    import json
    Path(__file__).with_name('beta_problem3_boxes.json').write_text(json.dumps(certified,indent=2)+'\n')

if __name__=='__main__':
    prove_positive(fmpq(2,3),fmpq(1))
    lo=fmpq(76859975974774087,10**17)
    hi=fmpq(76859975974774089,10**17)
    blo,clo=B_and_critical(arb(lo))
    bhi,chi=B_and_critical(arb(hi))
    assert clo<0 and chi>0
    print('CERTIFIED critical point in',lo,hi)
    B,_=B_and_critical(iv(lo,hi))
    assert B > arb(fmpq(59728368638449725,10**16))
    assert B < arb(fmpq(59728368638450275,10**16))
    print('CERTIFIED threshold enclosure =',B)
    for a,b in [(fmpq(4,3),fmpq(3,2))]:
        assert pg(0,arb(a))<0 and pg(0,arb(b))>0
    print('CERTIFIED unique digamma zero in (4/3,3/2)')
\end{filecontents*}
\documentclass[11pt,a4paper]{article}
\usepackage[T1]{fontenc}
\usepackage[utf8]{inputenc}
\usepackage{lmodern,microtype,amsmath,amssymb,amsthm,mathtools}
\usepackage[margin=25mm]{geometry}
\usepackage{booktabs,longtable,array,enumitem,xcolor,listings,fancyhdr}
\usepackage[colorlinks=true,linkcolor=blue!45!black,citecolor=blue!45!black,urlcolor=blue!45!black]{hyperref}
\usepackage{bookmark}
\hypersetup{pdftitle={Finite Criteria for Integration Stability and a Sharp Polygamma Parameter Problem},pdfauthor={Keio},pdfsubject={Research manuscript with proofs and reproducible verification programs}}
\setlength{\headheight}{14pt}
\setlength{\emergencystretch}{1em}
\pagestyle{fancy}\fancyhf{}
\fancyhead[L]{\small Keio}
\fancyhead[R]{\small Integration stability and polygamma parameters}
\fancyfoot[C]{\thepage}
\newtheorem{theorem}{Theorem}[section]
\newtheorem{proposition}[theorem]{Proposition}
\newtheorem{lemma}[theorem]{Lemma}
\newtheorem{corollary}[theorem]{Corollary}
\theoremstyle{definition}
\newtheorem{definition}[theorem]{Definition}
\newtheorem{example}[theorem]{Example}
\newtheorem{remark}[theorem]{Remark}
\newcommand{\C}{\mathbb C}
\newcommand{\Q}{\mathbb Q}
\newcommand{\Z}{\mathbb Z}
\newcommand{\N}{\mathbb N}
\DeclareMathOperator{\Stab}{Stab}
\DeclareMathOperator{\Hom}{Hom}
\DeclareMathOperator{\Gal}{Gal}
\DeclareMathOperator{\ord}{ord}
\DeclareMathOperator{\Tr}{Tr}
\lstset{language=Python,basicstyle=\ttfamily\scriptsize,breaklines=true,
columns=fullflexible,keepspaces=true,showstringspaces=false,
frame=single,rulecolor=\color{black!20},numbers=left,numberstyle=\tiny\color{black!50},
xleftmargin=1.5em,framexleftmargin=1em}
\title{Finite Criteria for Integration Stability\\
and a Sharp Polygamma Parameter Problem}
\author{Keio\thanks{Research development, independent internal mathematical audits,
and manuscript preparation with ChatGPT. No external peer review or
proof-assistant certification is claimed.}}
\date{2 November 2026}
\begin{document}
\maketitle
\begin{abstract}
We study three questions in symbolic integration and special functions.
First, finite monodromy representations describe the algebraic functions
that admit indefinitely repeated algebraic or elementary integration.
Affine homology controls algebraic primitives, while projective homology
controls elementary primitives. The proof combines the existing
Gavrilov--Pakovich moment theorem with finite group algebras, valuations,
and the geometry of divisors. Second, for a differential function field
of transcendence degree one, a function of field degree $d$ that is
D-finite has order at most $2d$. Normalized pole residues turn this into
a terminating existence test for creative telescoping in one primitive
extension. Third, we determine the complete admissible parameter interval
in Open Problem~3 of Qi--Lim--Nantomah (2025): it is
$[-1,\beta_*)$, where $\beta_*\approx5.972836863845$ is specified by a
unique, rigorously isolated critical point. The latter proof includes one
finite ball-arithmetic certificate. The sources and scope of each claim
are explicit. Neither unrestricted Chen--Kauers Problem~6 nor Problem~8
is claimed solved; this manuscript is not a report of one hundred
conjecture closures. Historical priority and a new foundational method
are not certified.
\end{abstract}

\section{Problems, scope, and verification}
This manuscript records the results that survived proof checks in a
research round drawing on the Verifiable Mathematical Reasoning Corpus,
at commit\par\noindent\texttt{ecd87de3348dcb706a76bb12fde722d7a348a012}\par\noindent
(2 November 2026), and on the primary references below. Existing corpus
claims are not recounted as new results. An infinite family of examples
is not counted as an infinite family of independently posed conjectures.

\begin{center}\small
\begin{tabular}{p{0.29\linewidth}p{0.62\linewidth}}
\toprule
Historical target & Result and scope \\
\midrule
Chen (2022), algebraic stability question &
A representation-theoretic classification for arbitrary algebraic
function fields, extending the earlier polynomial-inverse counterexamples.\\[4pt]
Chen--Kauers Problem 8 &
Classification for algebraic inputs, with both algebraic and elementary
primitive requirements. General transcendental elementary inputs remain
outside the theorem.\\[4pt]
Chen--Kauers Problem 6 &
A full existence test in one rational primitive extension with arbitrary
coupled algebraic poles. General Liouvillian and difference towers are
not resolved.\\[4pt]
Qi--Lim--Nantomah Open Problem 3 &
The original strict inequality is solved for all real parameters and
all positive arguments, including both parameter endpoints.\\
\bottomrule
\end{tabular}
\end{center}

The proof methods use established mathematical tools. In particular,
the moment theorem, Liouville's theorem, normalization of derivations,
Riemann--Roch, and rigorous ball arithmetic are credited to prior work.
A targeted search did not locate identical versions of the results
stated here, but this is not a certification of historical firstness.
Other problems from the same polygamma article already have public
solutions and are explicitly excluded from new-result claims.

The single source contains the verification programs as
\texttt{filecontents*} blocks. Ordinary compilation extracts them;
no shell escape is required. The programs are supplementary exact
checks except for the explicitly identified ball-arithmetic lemma,
whose finite computation is part of the proof. They do not constitute
a formal proof-assistant development.

\tableofcontents
\clearpage
% Fragment. Uses shared L,G,Y,S,A,D and names its projector e_{\mathrm{aff}}.
% Labels are prefixed aff:. Bibliography: Chen2022, GP2014, SSS2025.
\section{Algebraic primitive towers and affine homology}

Put $D=d/dx$. An algebraic function $f$ is \emph{algebraically stable} if there
are algebraic functions $f_0=f,f_1,f_2,\ldots$ with $Df_{j+1}=f_j$.
Chen asks for a description of this class in Section~5 of
\cite{Chen2022}. The condition is stronger than the elementary-integrability
condition in Chen--Kauers Problem~8: logarithms are not allowed in the present
primitive tower.

Let $\pi:Y\to\mathbb P^1_x$ be a finite connected Galois cover and write
\[
 L=\C(Y),\qquad G=\operatorname{Gal}(L/\C(x)),\qquad
 S=\pi^{-1}(\infty),\qquad A=\mathcal O(Y\setminus S).
\]
Thus $A$ is the integral closure of $\C[x]$ in $L$. Deck transformations act on
functions by $g\cdot q=q\circ g^{-1}$; they preserve $A$ and commute with $D$.
Let $H=H_1(Y\setminus S,\C)$. For $V\in\operatorname{Irr}(G)$ let $e_V\in\C[G]$ be the
central idempotent of its simple block, and define
\begin{equation}\label{aff:eq:good}
 e_{\mathrm{aff}}=\sum_{V\not\subset H}e_V.
\end{equation}
Here $V\not\subset H$ means that its multiplicity is zero; the projector
can include nontrivial representations as well as the trivial one.

\begin{theorem}[Structural classification]\label{aff:thm:main}
With the notation above,
\begin{align}
 A\cap\operatorname{Stab}(D,L)&=e_{\mathrm{aff}}A,\label{aff:eq:regular-class}\\
 \operatorname{Stab}(D,L)&=\bigcup_{N\ge0}D^N(e_{\mathrm{aff}}A).\label{aff:eq:full-class}
\end{align}
The union is increasing. Allowing primitives in an algebraic extension of $L$
does not enlarge either class. If the finite branch values have inertia groups
$I_1,\ldots,I_r$, then
\begin{equation}\label{aff:eq:character-formula}
 [H:V]=(r-1)\dim V-\sum_{i=1}^r\dim V^{I_i}
              +\begin{cases}1&V=\mathbf1,\\0&V\ne\mathbf1.\end{cases}
\end{equation}
In particular the trivial representation is always permitted.
\end{theorem}

This is a finite representation-theoretic description, rather than a condition
requiring infinitely many integrations. Its main external input is
\cite[Theorem~2.1]{GP2014}; that moment theorem is credited throughout. The
structural deduction is the identification of all cycle conditions with a
two-sided ideal, followed by pole removal. We do not assert that this is the
first formulation of the consequence in the literature.

\subsection{Primitives, moments, and descent}

\begin{lemma}\label{aff:lem:stability}
Let $F$ be a differential field with constants $\C$ and $Dx=1$.
An element $q\in F$ is stable if and only if every $x^nq$, $n\ge0$, has a
primitive in $F$. Moreover $q$ is stable if and only if $Dq$ is stable.
\end{lemma}
\begin{proof}
Repeated integration by parts proves the forward implication. Conversely,
if $DH_j=x^jq$ for $0\le j<k$, then
\[
 F_k=\frac1{(k-1)!}\sum_{j=0}^{k-1}(-1)^j\binom{k-1}{j}
                 x^{k-1-j}H_j
 \quad\text{satisfies}\quad D^kF_k=q.
\]
These arbitrarily deep primitives give a coherent tower. Indeed,
$\ker D^k=\C[x]_{<k}$, by induction using $Dx=1$; discrepancies between
chosen levels can therefore be corrected by polynomial primitives.
Differentiation preserves stability. If $Dq$ is stable, a stable primitive
of $Dq$ differs from $q$ by a constant, which is also stable.
\end{proof}

\begin{lemma}[Trace descent]\label{aff:lem:trace}
If $f\in L$ has an algebraic $N$th primitive, it has one in $L$.
Algebraic stability is consequently unchanged by finite algebraic extension.
\end{lemma}
\begin{proof}
The primitive $g$ belongs to some finite extension $E/L$. Unique extension of
derivations in characteristic zero implies that differentiation commutes with
field trace. Thus
\[
 D^N\left(\frac{\operatorname{Tr}_{E/L}(g)}{[E:L]}\right)=f.
\]
Apply this at every depth and use the polynomial-kernel correction from
Lemma~\ref{aff:lem:stability}.
\end{proof}

For $q\in A$, all differentials $x^nq\,dx$ have poles only in $S$. Such a
differential is exact in $L$ if and only if all its periods on $Y\setminus S$
vanish. Vanishing periods make path integration single-valued, and the
small-loop periods give zero residues. Local primitives then extend
meromorphically across $S$, producing a global meromorphic function on $Y$.
Consequently
\begin{equation}\label{aff:eq:periods}
 q\in A\cap\operatorname{Stab}(D,L)
 \quad\Longleftrightarrow\quad
 \int_\gamma x^nq\,dx=0
 \quad(n\ge0,\ \gamma\in H_1(Y\setminus S,\C)).
\end{equation}

\subsection{The moment theorem and the group-algebra ideal}

Choose an embedded star in the $x$-plane joining a regular value to each finite
branch value, and let $\Gamma$ be its inverse image. This graph is a deformation
retract of $Y\setminus S$. Orient edges from lifts of the center toward lifts
of branch values. For each branch-value label $i$, choose a reference edge and
write its translates as $e_{i,h}=h e_{i,1}$, $h\in G$.

We use the following specialization of the established moment theorem
\cite[Theorem~2.1 and the subsequent edge-count formula]{GP2014}.
If $\gamma$ is a graph cycle with edge coefficients $c_{i,h}$, then, for $q\in A$,
all moments $\int_\gamma x^nq\,dx$ vanish if and only if
\begin{equation}\label{aff:eq:gp}
 \sum_{h\in G}c_{i,h}\,q(hp(z))=0
 \quad\text{for every label }i.
\end{equation}
Here $p(z)$ is a reference inverse branch for the chosen arm. The hypotheses
hold because $q\,dx$ has no poles away from $S$. The same criterion holds for
complex cycles by the Cauchy-integral argument in the cited proof: the
generating function is a sum of Cauchy integrals whose densities are precisely
these branch sums with the complex edge coefficients. Vanishing moments make
that function zero near infinity, and its jump across each arm forces the
corresponding density to vanish. The reverse implication follows by summing
the arm integrals. Thus the complex-cycle equivalence does not rely on
extending an integral-cycle equivalence by linearity alone.

Define
\[
 u_{\gamma,i}=\sum_{h\in G}c_{i,h}h^{-1}\in\C[G],\qquad
 U=\operatorname{span}_{\C}\{u_{\gamma,i}\},\qquad J=\C[G]U.
\]
The inverse in this formula is essential: $q(hp)=(h^{-1}\cdot q)(p)$.
Translating a cycle by $a\in G$ gives
\begin{equation}\label{aff:eq:right-ideal}
 u_{a\gamma,i}=u_{\gamma,i}a^{-1}.
\end{equation}
Thus $U$ is a right ideal and $J$ is two-sided. Conditions \eqref{aff:eq:gp} are
exactly $Uq=0$. Applying $a\in G$ to any identity $uq=0$ gives $(au)q=0$;
therefore
\begin{equation}\label{aff:eq:ideal-condition}
 Uq=0\quad\Longleftrightarrow\quad Jq=0.
\end{equation}

\begin{lemma}\label{aff:lem:blocks}
The ideal $J$ is the sum of the simple matrix blocks of $\C[G]$ indexed by
the irreducible types occurring in $H$.
\end{lemma}
\begin{proof}
Every edge orbit is regular, so
$C_1(\Gamma,\C)=\bigoplus_{i=1}^r\C[G]$. As a graph has no two-chains,
$H=\ker\partial$ embeds in this direct sum. Its coordinate projections have
exactly its collective irreducible support: each image is a quotient of the
semisimple module $H$, and a nonzero isotypic component must have a nonzero
coordinate projection because the embedding is injective.

Inversion in $\C[G]$ changes left modules into dual right modules. Hence the
right-ideal support of $U$ is dual to the support just described. A two-sided
ideal generated by a nonzero part of a simple matrix block contains that
entire block and no block absent from the generators. Finally $H$ has an
integral, hence real, character; its support is closed under duality. This
identifies the blocks of $J$ as claimed.
\end{proof}

The annihilator of these blocks on $L$ is $e_{\mathrm{aff}}L$, and its intersection with
$A$ is $e_{\mathrm{aff}}A$. Equations \eqref{aff:eq:periods}--\eqref{aff:eq:ideal-condition}
and Lemma~\ref{aff:lem:blocks} prove \eqref{aff:eq:regular-class}.

The character formula follows from the connected-graph sequence
\[
 0\longrightarrow H\longrightarrow C_1(\Gamma)
  \longrightarrow C_0(\Gamma)\longrightarrow\C\longrightarrow0.
\]
The center vertices form a regular orbit; those above branch value $i$ form
$G/I_i$. Thus
\[
 [H]=(r-1)[\C[G]]-\sum_i[\C[G/I_i]]+[\mathbf1].
\]
Taking multiplicities and using Frobenius reciprocity proves
\eqref{aff:eq:character-formula}. It also covers the trivial degree-one cover,
where $r=0$ and $G$ is trivial.

\subsection{Removing finite poles and a finite decision test}

For nonzero $f\in L$ and $P\notin S$, set
\[
 e_P=\operatorname{ord}_P(x-x(P)),\qquad m_P=\max(0,-\operatorname{ord}_Pf),\qquad
 N_f=\max_{P\notin S}\left\lceil\frac{m_P}{e_P}\right\rceil,
\]
with an empty maximum equal to zero. Only the finitely many poles of $f$
contribute.

\begin{lemma}[Pole removal]\label{aff:lem:poles}
If $D^{N_f}q=f$ with $q\in L$, then $q\in A$.
\end{lemma}
\begin{proof}
At a finite point choose a local parameter $t$ with
$x-x(P)=t^{e_P}$. Then
$D=(e_Pt^{e_P-1})^{-1}d/dt$. A pole of order $k>0$ of $q$ therefore gives
a pole of order exactly $k+N_fe_P$ of $f$, with a nonzero leading
coefficient. But $m_P\le N_fe_P$, a contradiction.
\end{proof}

If $f$ is stable, trace descent supplies such a $q$. It is itself stable:
given $j$, choose $h$ with $D^{N_f+j}h=f$. Then
$D^jh-q\in\ker D^{N_f}=\C[x]_{<N_f}$, and subtracting a polynomial
$j$th primitive from $h$ gives a $j$th primitive of $q$.
Lemma~\ref{aff:lem:poles} and \eqref{aff:eq:regular-class} imply $q\in e_{\mathrm{aff}}A$.
The converse in \eqref{aff:eq:full-class} follows because derivatives preserve
stability. A primitive of a stable regular element has no finite pole, again
by the local pole calculation; hence it belongs to $e_{\mathrm{aff}}A$. This proves that
the union is increasing and completes Theorem~\ref{aff:thm:main}.

\begin{theorem}[Finite Riemann--Roch certificate]\label{aff:thm:finite}
For nonzero $f\in L$, set $N=N_f$. At $P\in S$ write $e_P=-\operatorname{ord}_Px$ and put
\[
 B_P=\max\{Ne_P-\operatorname{ord}_Pf,(N-1)e_P,0\},\qquad
 E_f=\sum_{P\in S}B_PP.
\]
Then $f$ is algebraically stable if and only if the finite-dimensional linear
system
\begin{equation}\label{aff:eq:finite-system}
 q\in\mathcal L(E_f),\qquad (1-e_{\mathrm{aff}})q=0,\qquad D^Nq=f
\end{equation}
has a solution. The zero function is stable separately.
\end{theorem}
\begin{proof}
It remains only to bound the poles at infinity of a primitive $q$. Choose
a local parameter $z$ with $x=z^{-e_P}$, so
$D=-e_P^{-1}z^{e_P+1}d/dz$. On a Laurent monomial $z^{-k}$, $D^N$ shifts
the exponent by $Ne_P$ and has nonzero coefficient unless
$k\in\{0,e_P,\ldots,(N-1)e_P\}$ (the exceptional set is empty when $N=0$).
Distinct exponents remain distinct.
If a leading pole is not exceptional, its order is
$k=Ne_P-\operatorname{ord}_Pf$; if it is exceptional, $k\le(N-1)e_P$.
This proves the bound. Necessity and sufficiency now follow from
Theorem~\ref{aff:thm:main} and Lemma~\ref{aff:lem:poles}.
\end{proof}

For an effectively presented algebraic function, a finite procedure is thus:
construct a finite Galois normal closure and its group action; compute the
finite inertia data and \eqref{aff:eq:character-formula}; compute the divisor of
$f$, the bound $E_f$, and a basis of $\mathcal L(E_f)$; then solve
\eqref{aff:eq:finite-system} by linear algebra. Equality in the function field
reduces to coefficient equality in a finite algebraic basis over $\C(x)$.
This is a mathematical finite procedure with the usual effective curve and
constant-field operations supplied. It is not a computability claim for
arbitrary unnamed complex constants.

For a non-Galois field $K=\C(x,f)$, apply the procedure in a normal closure
$L$. A successful certificate gives algebraic stability; conversely stability
in the algebraic closure descends to $K$ and hence to $L$ by
Lemma~\ref{aff:lem:trace}. No extra Galois hypothesis on the input is required.

\subsection{Examples and consistency checks}

\paragraph{Cyclic powers.}
For $x=t^m$ every irreducible type is absent from affine homology. Therefore
\eqref{aff:eq:full-class} gives
\[
 \operatorname{Stab}(D,\C(t))=\C[x]\oplus
       \bigoplus_{i=1}^{m-1}\C[x,x^{-1}]t^i.
\]
Indeed $D^Nt^k$ is a constant multiple of $t^{k-Nm}$; the multiple vanishes
exactly when a polynomial in $x$ is differentiated too often. This agrees
with the existing classification in \cite[Proposition~5.3(iii)]{SSS2025}.

\paragraph{A biquadratic cover.}
Let $s^2=x$ and $t^2=x-1$. The group is $C_2\times C_2$ and the two finite
inertia groups flip $s$ and $t$, respectively. Formula
\eqref{aff:eq:character-formula} assigns multiplicity zero to the trivial
character and to the characters of $s$ and $t$, and multiplicity one to
their product. Thus $s+t$ is stable, while $st=\sqrt{x(x-1)}$ is not
algebraically stable. Explicitly,
\[
 D\left(\tfrac23s^3+\tfrac23t^3\right)=s+t,
\]
and continued integration stays algebraic term by term. The latter product
is nevertheless elementary-stable, illustrating the distinction in scope.

\paragraph{Hyperelliptic fields.}
For $y^2=P(x)$ with $P$ squarefree of degree $d\ge2$, the sign representation
has affine-homology multiplicity $d-1$ and is forbidden. The only permitted
regular functions are $\C[x]$. Since their derivatives remain polynomials,
\[
 \operatorname{Stab}(D,\C(x,y))=\C[x].
\]
This includes positive-genus examples and is a full statement about the
function field, not merely its regular coordinate ring.

\paragraph{The cubic normal closure.}
For $x=u^3+u$, the normal closure has group $S_3$ and two transposition inertia
groups. The multiplicities of the trivial, standard and sign representations
are respectively $0,0,1$. The standard component is allowed, in agreement
with the polynomial primitive tower for the cubic inverse. This permits
nonabelian examples excluded by the specific radical-sum description proposed
in \cite{Chen2022}.

\paragraph{A Laurent cover.}
Let $x=t^m+t^{-m}$, $m\ge2$. The Galois group is dihedral, the covering curve
is $\mathbb P^1_t$, and $S=\{0,\infty\}$. Affine homology is one-dimensional:
rotations act trivially and inversion $t\mapsto t^{-1}$ acts by minus one.
Thus, with $A=\C[t,t^{-1}]$,
\[
 A\cap\operatorname{Stab}(D,\C(t))
   =\C[x]+\operatorname{span}_{\C}\{t^k:k\in\Z,\ m\nmid k\}.
\]
The invariant symmetric terms $t^{jm}+t^{-jm}$ are polynomials in $x$.
For $m\nmid k$ the rule
\[
 \int t^k\,dx=
  \frac{m}{k+m}t^{k+m}-\frac{m}{k-m}t^{k-m}
\]
never reaches a zero denominator under iteration. The forbidden rotation-
invariant antisymmetric part produces logarithmic obstructions. In particular
$t^m+t^{-m}=x$ is stable, while $t^m-t^{-m}$ is not algebraically stable.

\begin{corollary}[The exact abelian scope of the proposed class]
\label{aff:cor:abelian-scope}
Inside the maximal abelian extension of $\C(x)$, the algebraically stable
elements are exactly the finite sums
\[
 \sum_i p_i(x)(x-c_i)^{\alpha_i},\qquad
 p_i\in\C[x],\quad c_i\in\C,\quad
 \alpha_i\in\Q\setminus\{-1,-2,\ldots\}.
\]
Thus Chen's proposed description is correct with this abelian restriction.
\end{corollary}
\begin{proof}
Work in a finite abelian Galois extension containing the element. For a
nontrivial character $\chi$, let $q_\chi$ count the finite inertia groups on
which it is nontrivial. Formula~\eqref{aff:eq:character-formula} gives
$[H:\chi]=q_\chi-1$. Hence a permitted nontrivial character is ramified at
exactly one finite value $c$. The quotient by $\ker\chi$ is cyclic and is
ramified only at $c$ and infinity. Kummer theory makes its field
$\C((x-c)^{1/m})$: in a Kummer equation $y^m=R(x)$, every other finite
valuation of $R$ is divisible by $m$ and can be removed by an $m$th power;
the remaining exponent at $c$ is coprime to $m$.

The regular $\chi$-component consists of polynomial multiples of
$(x-c)^{a/m}$ with $0<a<m$, and its derivative closure consists of finite
sums of noninteger rational powers centered at $c$. The trivial component
contributes $\C[x]$. Theorem~\ref{aff:thm:main} proves necessity. Conversely,
each displayed term admits algebraic primitives at every depth, and a finite
collection of the indicated radicals lies in a finite abelian Kummer
compositum. Trace descent ensures that allowing primitives outside that
compositum does not change the necessity argument.
\end{proof}

\subsection{Relation to earlier work and remaining limits}

The all-cycle moment criterion is already provided by Gavrilov--Pakovich
\cite{GP2014}. The present deduction organizes it into central representation
blocks and combines that organization with an explicit finite-pole and
Riemann--Roch reduction. Sahu--Seiss--Srinivasan \cite{SSS2025} study finite
elementary integration and classify several specific differential fields,
including the cyclic radical case checked above. Their stated results do not
give the arbitrary finite-cover representation criterion used here.

The conclusion concerns Chen's algebraic-stability classification question.
It does not establish a classification of all elementary-stable elementary
functions, or even invoke such a classification. No assertion is made that
every permitted representation is generated by genus-zero polynomial
subcovers; that stronger compositional reformulation would need another proof.
The characterization is one result related to a pre-existing question, not
many separate conjecture resolutions. A literature search cannot certify
absence of prior formulations, so no historical firstness claim is included.


\section{Elementary stability of arbitrary algebraic inputs}
\label{sec:elementary-monodromy}

We now allow elementary primitives, while retaining algebraic inputs.  This
distinction changes the controlling representation from affine to projective
homology.  The result below does not classify transcendental elementary
functions and hence does not settle the full scope of Problem~8
of~\cite{CK2016}.

Retain the notation $L=\C(Y)$, $G$, $x$, $S=x^{-1}(\infty)$,
$A=\mathcal O(Y\setminus S)$, and $D=d/dx$ for a connected finite Galois
cover $x:Y\to\mathbf P^1$.  A function is \emph{elementary-stable} if it has
an elementary primitive of every order.  Let
\[
 H_{\mathrm{proj}}=H_1(Y,\C),\qquad
 e_{\mathrm{proj}}
   =\sum_{\rho\notin\operatorname{supp}H_{\mathrm{proj}}}e_\rho,
\]
where $e_\rho$ is the central primitive idempotent of $\C[G]$ associated
with an irreducible representation $\rho$.  The set of summands is invariant
under Galois conjugacy of characters, because $H_1(Y,\Q)$ is a rational
representation.  Thus $e_{\mathrm{proj}}\in\Q[G]$.  It is also fixed by
the involution $g\mapsto g^{-1}$, since projective homology is self-dual.
We use the action $(g f)(P)=f(g^{-1}P)$ on functions, and the corresponding
action on differentials.  Divisors and homology carry the deck action.

\begin{theorem}[Projective-homology criterion]
\label{thm:elementary-projective}
The elementary-stable functions in $L$ are precisely
\[
 \operatorname{Stab}_{\mathrm{elem}}(L)=e_{\mathrm{proj}}L.
\]
If $K=L^J$ is an intermediate field, then
\[
 \operatorname{Stab}_{\mathrm{elem}}(K)
   =e_{\mathrm{proj}}L\cap K=e_{\mathrm{proj}}K.
\]
In particular, the stable set is a $\C(x)$-vector space, invariant under $D$.
\end{theorem}

Elementary stability is equivalent to elementary integrability of every
moment $x^n f$, $n\geq0$.  In one direction, integrate by parts using an
$(n+1)$st primitive of $f$.  In the other direction, if $DI_j=x^j f$ for
$0\leq j<N$, then
\[
 \frac1{(N-1)!}\sum_{j=0}^{N-1}
    (-1)^j\binom{N-1}{j}x^{N-1-j}I_j
\]
is an $N$th primitive of $f$.  Finite sums and products by polynomials
therefore preserve stability.  Initially independent choices for the
different orders can be adjusted to a compatible tower.  Indeed,
$\ker D^N=\C[x]_{<N}$ in an elementary closure with constant field $\C$;
the difference between a chosen $N$th primitive and the derivative of an
$(N+1)$st primitive is such a polynomial and can be removed by a
polynomial correction.

We first prove the sufficient condition in a stronger form.

\begin{lemma}[Differentials supported on absent representations]
\label{lem:good-differential}
Every meromorphic differential $\omega$ on $Y$ satisfying
$e_{\mathrm{proj}}\omega=\omega$ has an elementary primitive.
\end{lemma}

\begin{proof}
Write $e=e_{\mathrm{proj}}$, and choose a finite $G$-invariant set $T$
containing the poles of $\omega$.  Its residue vector $r$ belongs to
$\operatorname{Div}^0_T\otimes_{\Z}\C$, by the residue theorem, and
$er=r$ by equivariance.  Here $\operatorname{Div}^0_T$ is the lattice of
degree-zero divisors supported on $T$.

Choose a positive integer $M$ such that $Me\in\Z[G]$.  This group-ring
element induces an endomorphism of $\operatorname{Jac}(Y)$, with zero
action on $H_1(\operatorname{Jac}(Y),\Z)\simeq H_1(Y,\Z)$.  It is
therefore the zero endomorphism: its complex-linear lift to the universal
cover vanishes on the full period lattice.  Consequently, for every
$B\in\operatorname{Div}^0_T$, the integral divisor $MeB$ has trivial
Jacobian class and is principal.

Take a lattice basis $B_1,\ldots,B_a$ of $\operatorname{Div}^0_T$ and
write $r=\sum_j c_j eB_j$.  Choose $u_j\in L^*$ with
$\operatorname{div}(u_j)=MeB_j$.  The logarithmic differential
\[
 \eta=\frac1M\sum_j c_j\frac{du_j}{u_j}
\]
has residue vector $r$.  The differential $e\eta$ remains a finite
$\C$-linear combination of logarithmic derivatives, is fixed by $e$, and
has residue vector $er=r$.  Thus $\theta=\omega-e\eta$ has zero
residues and is fixed by $e$.

The periods of a differential with zero residues define a functional on
$H_1(Y,\C)$.  The period map is $G$-equivariant, and its target has zero
$e$-part.  Hence every period of $\theta$ is zero.  Path integration now
gives a single-valued primitive; zero residues ensure that its local
Laurent expansions are meromorphic at the poles.  Therefore $\theta=da$
for some $a\in L$.  The identity $\omega=da+e\eta$ is an elementary
primitive certificate.
\end{proof}

To prove necessity, we require periods that annihilate every allowed
logarithmic term.  Put $H_{\mathrm{aff}}=H_1(Y\setminus S,\C)$ and
$U_S=A^*/\C^*$.  The divisor map embeds $U_S$ into
$\operatorname{Div}^0_S$, so $U_S$ is free abelian of finite rank.  Define
\[
 \ell_u(\gamma)=\frac1{2\pi i}\int_\gamma\frac{du}{u},\qquad
 W_S=\bigcap_{u\in A^*}\ker\ell_u\subseteq H_{\mathrm{aff}}.
\]
The functionals $\ell_u$ are integer-valued on integral homology.
Consequently $W_S$ is a rational $G$-submodule; its irreducible support
is closed under duality.

\begin{lemma}[Unit-adjusted periods]
\label{lem:unit-periods}
\begin{enumerate}
\item Every elementary-integrable meromorphic differential with poles only
in $S$ has zero periods on $W_S$.
\item The natural map $W_S\to H_{\mathrm{proj}}$ is surjective.
\end{enumerate}
\end{lemma}

\begin{proof}
For the first assertion, Liouville's theorem
(as stated in~\cite[Theorem~3.2]{Chen2022}) and trace from a finite
algebraic extension give
\[
 \omega=da+\sum_i c_i\frac{db_i}{b_i},\qquad
 a\in L,\quad b_i\in L^*,\quad c_i\in\C.
\]
Trace descends the exact term, while logarithmic derivatives descend by
norms.  The function $a$ has no pole outside $S$: such a pole would give
$da$ a pole of order at least two, whereas logarithmic derivatives have
at most simple poles.

Choose a $\Q$-linearly independent basis $\beta_1,\ldots,\beta_b$ of
the $\Q$-span of the $c_i$.  Clearing denominators and taking products
of the $b_i$ rewrites the logarithmic sum as
\[
 \frac1N\sum_{j=1}^b\beta_j\frac{dv_j}{v_j},\qquad v_j\in L^*.
\]
At $P\notin S$, its residue is zero.  The equality
$\sum_j\beta_j\operatorname{ord}_P(v_j)=0$ and $\Q$-linear
independence force every order to vanish.  Thus all $v_j$ are $S$-units,
and integration over $W_S$ kills both the exact and logarithmic terms.

For the second assertion, the boundary-loop sequence is
\[
 0\longrightarrow B_S\longrightarrow H_{\mathrm{aff}}
   \longrightarrow H_{\mathrm{proj}}\longrightarrow0,
\]
where $B_S$ is generated by small loops about the points of $S$, subject
to their single sum relation.  Choose a basis $u_1,\ldots,u_b$ of $U_S$.
The map
\[
 B_S\longrightarrow\C^b,
 \qquad \gamma\longmapsto
       (\ell_{u_1}(\gamma),\ldots,\ell_{u_b}(\gamma))
\]
is surjective.  Indeed, its rows are the integer vectors
$\operatorname{div}(u_j)$, which are independent over $\Q$ and hence
over $\C$.  Any lift of a projective homology class can therefore be
adjusted by a boundary class to kill all unit windings.  The adjusted
lift belongs to $W_S$.
\end{proof}

\begin{proof}[Proof of Theorem~\ref{thm:elementary-projective}]
If $f=e_{\mathrm{proj}}f$, then every $x^n f\,dx$ is fixed by
$e_{\mathrm{proj}}$, because $x$ and $dx$ are $G$-invariant.
Lemma~\ref{lem:good-differential} and the moments criterion prove
elementary stability.

Conversely, first let $q\in A$ be elementary-stable.  By
Lemma~\ref{lem:unit-periods},
\begin{equation}
 \int_\gamma x^n q\,dx=0
 \qquad(\gamma\in W_S,\ n\geq0).
 \label{eq:unit-all-moments}
\end{equation}
Use the constellation graph of $x$: the inverse image of a star joining
an ordinary value to all finite branch values.  This graph is a
deformation retract of $Y\setminus S$.  Each family of edges above one
star arm is a regular $G$-set.

The moment theorem of Gavrilov--Pakovich~\cite[Theorem~2.1 and
formula~(15)]{GP2014} turns~\eqref{eq:unit-all-moments} into branch-sum
identities on every arm.  Its Cauchy-integral proof applies also to
complex cycles: the generating function is a sum of Cauchy integrals
whose densities are the branch sums with the complex edge coefficients.
Zero moments give the zero generating function near infinity, and its
jump across each arm forces that density to vanish.  Thus the use of
the complex subspace $W_S$ is justified directly by the proof, without
requiring a basis consisting of individual closed loops.

Label the edges over an arm by $e_h=h e_1$.  If a cycle has coefficients
$c_h$ on that arm, the associated identity is
\[
 u q=0,
 \qquad u=\sum_h c_h h^{-1}\in\C[G].
\]
Let $U$ be the sum of the images of all arm projections of $W_S$,
followed by this inversion.  Deck translation by $a$ sends $u$ to
$ua^{-1}$, so $U$ is a right ideal.  Applying a deck automorphism to
$uq=0$ gives $(au)q=0$; consequently $Jq=0$ for the two-sided ideal
$J=\C[G]U$.

The arm projections together embed $W_S$ into a direct sum of regular
representations.  Their combined irreducible support is exactly that
of $W_S$.  Inversion dualizes the corresponding blocks, but this does
not change the support because $W_S$ is rational.  By semisimplicity,
$J$ therefore contains precisely the blocks occurring in $W_S$.
Lemma~\ref{lem:unit-periods} shows that these include every block of
$H_{\mathrm{proj}}$.  Hence $(1-e_{\mathrm{proj}})q=0$.

For arbitrary stable $f\in L$, choose a nonzero $R\in\C[x]$ such
that $q=R(x)f\in A$.  This is possible by assigning sufficiently high
zeros at the finitely many finite $x$-values below its poles.  Polynomial
multiplication preserves stability, so the preceding argument gives
$q=e_{\mathrm{proj}}q$.  Division by $R(x)$ yields
$f=e_{\mathrm{proj}}f$.

Finally, a finite algebraic extension does not alter elementary
integrability, by Liouville trace and norm descent.  Applying this to
all moments proves the intermediate-field assertion.  Centrality of
$e_{\mathrm{proj}}$ makes it preserve the $J$-fixed field, and its
commutation with $D$ and with $\C(x)$ proves the remaining assertions.
\end{proof}

The projector is determined by finite monodromy data.  If the finite
branch values have inertia groups $I_1,\ldots,I_r$ and the inertia
group at infinity is $I_\infty$, the graph chain complex, followed by
removal of the boundary-loop module, gives
\[
 [H_{\mathrm{proj}}]
  =(r-1)[\C[G]]-\sum_{i=1}^r[\C[G/I_i]]
       -[\C[G/I_\infty]]+2[\mathbf1].
\]
Thus the multiplicity of $\rho$ is
\[
 (r-1)\dim\rho-\sum_{i=1}^r\dim\rho^{I_i}
       -\dim\rho^{I_\infty}+2\delta_{\rho,\mathbf1}.
\]
This is a structural finite-data criterion, with no assumption of
computability for arbitrary complex coefficients.

\begin{example}
If $Y$ has genus zero, $e_{\mathrm{proj}}=1$ and every element of $L$
is elementary-stable.  This agrees with rational parametrization.
If $L=\C(x,y)$ with $y^2=P(x)$, where $P$ is squarefree of degree at
least three, then $Y$ has positive genus and the hyperelliptic
involution acts as $-1$ on $H_1(Y,\C)$.  The good projector is the
averaging projector, so the elementary-stable elements of this field
are exactly $\C(x)$.
\end{example}

\begin{example}[A stable subspace on a positive-genus Galois curve]
Consider
\[
 L=\C(x,\sqrt{x},\sqrt{x-1},\sqrt{x-2}),\qquad G\simeq(C_2)^3.
\]
The inertia generators at $0,1,2$ are the three coordinate vectors,
and the generator at infinity is their sum. Riemann--Hurwitz gives
$2g(Y)-2=8(-2)+4\cdot4=0$, so $g(Y)=1$.
For a nontrivial character supported on $j$ coordinates, the number
of branch points on which it is nontrivial is $j+(j\bmod2)$.
The projective multiplicity is therefore zero when $j=1$ or $2$,
and two when $j=3$. It follows that the elementary-stable space has
$\C(x)$-dimension seven, with basis
\[
 1,\ \sqrt{x},\ \sqrt{x-1},\ \sqrt{x-2},\
 \sqrt{x(x-1)},\ \sqrt{x(x-2)},\ \sqrt{(x-1)(x-2)}.
\]
The remaining basis element $\sqrt{x(x-1)(x-2)}$ is not elementary-stable.
Thus the criterion also identifies proper stable subspaces inside
positive-genus Galois function fields.
\end{example}

\begin{corollary}[The abelian case]
\label{cor:elementary-abelian}
Inside the maximal abelian algebraic extension of $\C(x)$, the
elementary-stable elements are exactly the finite sums
\[
 \sum_{j=1}^m a_j(x)\,u_j(x)^{r_j},
 \qquad a_j\in\C(x),\quad r_j\in\Q,
\]
where each $u_j$ is a nonconstant fractional linear function.
\end{corollary}

\begin{proof}
Work first in a finite abelian Galois extension.  For a nontrivial
one-dimensional character $\chi$, the displayed multiplicity formula
simplifies to $b_\chi-2$, where $b_\chi$ counts all branch points,
including infinity, whose inertia has nontrivial image under $\chi$.
Thus $\chi$ is absent from projective homology exactly when
$b_\chi=2$.

The cyclic quotient cover corresponding to $\ker\chi$ consequently
has exactly two branch points.  Such a connected cover is totally
ramified at both points and is obtained by taking a root of a
fractional linear function carrying those points to $0$ and $\infty$.
Its character eigenspaces are one-dimensional over $\C(x)$ and are
generated by rational powers of that fractional linear function.
Character decomposition gives the asserted form.  Conversely, each
displayed summand belongs to a cyclic genus-zero function field and
is elementary-stable; finite sums preserve stability.  Taking the
union over finite abelian extensions completes the proof.
\end{proof}

The proof combines Liouville theory, the Jacobian, and the published
moment theorem.  The classification above is a conclusion for
algebraic inputs, not an assertion that the underlying methods or all
of their consequences are historically new.

\section{A finite existence criterion in a single primitive extension}
\label{int:section}

This section gives a terminating criterion for arbitrary rational inputs in
one primitive extension.  It does not settle Chen--Kauers Problem~6 in its
full scope: nested Liouvillian towers and difference fields are not covered.
The proof uses classical normalization and algebraic-curve theory.  Its
historical novelty is not established by the literature checks reported here.

\subsection{A degree bound for differential finiteness}

Let $(k,d)$ be a differential field of characteristic zero, and let
$(K,\delta)$ be a finitely generated differential extension of transcendence
degree one.  Thus $\delta|_k=d$.  We call $r\in K$ \emph{$D$-finite over $k$}
if $\operatorname{span}_k\{\delta^j r:j\geq0\}$ is finite dimensional.
The algebraic closure of $k$ inside a function field is called its
\emph{relative algebraic constant field}; this is different from its field
of $\delta$-constants.

\begin{lemma}[Finite generation of the relevant differential subfield]
\label{int:stabilize}
If $r$ is transcendental over $k$, put $d_0=[K:k(r)]$ and
\[
 E_i=k(r,\delta r,\ldots,\delta^i r).
\]
There are at most $\lfloor\log_2 d_0\rfloor$ strict inclusions in this
chain.  At the first equality $E_i=E_{i+1}$, the field $E_i$ is
$\delta$-stable and equals $E=k\langle r\rangle$.
\end{lemma}
\begin{proof}
Every strict inclusion multiplies $[E_i:k(r)]$ by an integer at least two,
while $[E_i:k(r)]\leq d_0$.  At an equality, differentiation sends each
field generator of $E_i$ into $E_i$; the quotient rule then gives
$\delta E_i\subseteq E_i$.  Hence a plateau cannot be followed by a
strict inclusion.
\end{proof}

Let $X$ be the nonsingular complete curve over $k$ with function field
$E=k\langle r\rangle$, and let $D=(r)_\infty$ be the pole divisor of $r$.
We allow a finite relative algebraic constant extension of $k$, so $X$
need not be geometrically connected over $k$.

\begin{theorem}[Curve criterion]
\label{int:curve}
For $r$ transcendental over $k$, the following conditions are equivalent:
\begin{enumerate}
\item $r$ is $D$-finite over $k$;
\item $\delta\mathcal O_{X,P}\subseteq\mathcal O_{X,P}$ at every closed
point $P$, and $\delta\mathfrak m_P\subseteq\mathfrak m_P$ at every pole
of $r$;
\item $\delta\mathcal L(D)\subseteq\mathcal L(D)$, where
$\mathcal L(D)$ is the Riemann--Roch space of $D$.
\end{enumerate}
When these conditions hold, the minimal scalar differential order of $r$
is at most $\dim_k\mathcal L(D)$.
\end{theorem}
\begin{proof}
Suppose first that the orbit space $W$ of $r$ is finite dimensional;
adjoining $1$ to $W$ is harmless.  The ring $A=k[W]$ is a finitely
generated $\delta$-stable domain with fraction field $E$.
Seidenberg's normalization theorem in characteristic zero~\cite{Seidenberg1966} implies that
its integral closure $B$ in $E$ is also $\delta$-stable.
Let $S$ be the union of the pole supports of a basis of $W$ on $X$.
The normalized affine curve $\operatorname{Spec}B$ is $X\setminus S$.
Localization therefore proves regularity of $\delta$ away from $S$.

Fix $P\in S$ and a local parameter $z$.  The completed local field has
a coefficient field which is finite algebraic over $k$: lift a separating
primitive element of the residue field by Hensel's lemma.  The unique
extension of $d$ preserves this coefficient field.  Thus the derivation
on Laurent series consists of coefficient differentiation and
$(\delta z)\,\partial_z$.
Choose $w\in W$ with maximal pole order $m>0$ at $P$ and write
$w=c z^{-m}+\cdots$, with $c\ne0$.  If
$v_P(\delta z)=\ell\leq0$, and $b$ is the leading coefficient of
$\delta z$, then
\[
 \delta w=-mcb\,z^{-m-1+\ell}+\text{higher terms}.
\]
Coefficient differentiation has exponent at least $-m$, so this leading
term cannot cancel.  It contradicts maximality of $m$, since
$\delta w\in W$.  Hence $v_P(\delta z)\geq1$: the derivation is regular
and tangent to the maximal ideal at $P$.  This proves condition~2.

Under condition~2, differentiation preserves regularity outside the poles
of $r$.  At a pole, write $h=z^n u$ with $u$ a local unit.  Tangency and
the product rule give $v_P(\delta h)\geq v_P(h)$.  Every defining
valuation bound of $\mathcal L(D)$ is consequently preserved, proving
condition~3.  Finally, this finite-dimensional space contains $r$, so
the first $\dim_k\mathcal L(D)+1$ derivatives are linearly dependent.
\end{proof}

\begin{corollary}[Explicit stopping bound]
\label{int:bound}
For transcendental $r\in K$, if $r$ is $D$-finite over $k$, its minimal
order is at most
\[
                    2[K:k(r)].                         \tag{I.1}
\]
If $k$ is relatively algebraically closed in $K$, the sharper bound
$[K:k(r)]+1$ holds.  If $r$ is algebraic over $k$, it is $D$-finite of
order at most $[k(r):k]$.
\end{corollary}
\begin{proof}
Let $c_E$ be the relative algebraic constant field of $E$,
$q=[c_E:k]$, and $e=[E:k(r)]$.  The pole divisor of $r$ has degree
$e$ over $k$, or $e/q$ over $c_E$.  Therefore
\[
 \dim_k\mathcal L(D)
 =q\dim_{c_E}\mathcal L(D)
 \leq q(e/q+1)=e+q\leq2e\leq2[K:k(r)].
\]
Here we used the elementary Riemann--Roch estimate
$\dim\mathcal L(D)\leq\deg D+1$ and $q\leq e$.
In the regular case $q=1$.  For algebraic $r$, differentiating its
separable minimal polynomial shows that $k(r)$ is $\delta$-stable, and
finite-dimensional linear algebra applies directly.
\end{proof}

Consequently $D$-finiteness in a supplied transcendence-degree-one
differential function field can be decided by a finite computation.
For transcendental $r$, compute $d_0=[K:k(r)]$ and test
\[
                  r,\delta r,\ldots,\delta^{2d_0}r
                  \quad\text{for dependence over }k.       \tag{I.2}
\]
The earliest dependence gives the minimal operator; absence of a
dependence proves nonexistence.  These are effective algebraic operations
over a computably presented characteristic-zero field with factorization.
For example, represent $K$ over $k(y)$, clear denominators in the
coordinates of the tested elements, and compare coefficients of $y$.
The field degree in (I.1) is not the degree of an unreduced numerator in
an arbitrary algebraic presentation.  Normalization and Riemann--Roch
are used in the proof, but need not be implemented to execute (I.2).

\subsection{Normalized residues at arbitrary algebraic poles}

Let $C$ be algebraically closed of characteristic zero, let $k=C(x)$,
and choose $a,b\in k(y)$ satisfying $a_x=b_y$, with $a$ having no
rational antiderivative with respect to $y$ in $k(y)$.
Let $t$ be transcendental and set
\[
 F=C(x,y,t),\qquad D_x=\partial_x+b\partial_t,\qquad
 D_y=\partial_y+a\partial_t.                            \tag{I.3}
\]
The compatibility condition makes these derivations commute.  This
includes primitives of arbitrary rational closed one-forms, such as
$\log(x+y)$ and $\log(x^2+y^2)$.
We seek nonzero $P\in k[D_x]$ and $g\in F$ such that $Pf=D_yg$.
All assertions concern certificates in this specified field.

For an arbitrary $f\in F$, choose a finite normal algebraic extension
$K/k(y)$ splitting its denominator as a polynomial in $t$.
Both coefficient derivations extend uniquely to $K$ and commute.
At a root $\alpha\in K$, put
\[
 p=t-\alpha,\qquad \eta=a-\partial_y\alpha,\qquad
 \nu=\partial_x\alpha-b,\qquad v=\frac\nu\eta,\qquad
 \delta_\alpha=\partial_x+v\partial_y.                  \tag{I.4}
\]
Here $\eta\ne0$.  Indeed, an algebraic primitive of $a$ would yield,
by normalized trace to $k(y)$, a rational primitive of $a$, contrary to
the hypothesis.

Write the principal part of $f$ at $p$ as
$\sum_{j=1}^m c_jp^{-j}$, and define the $k$-linear operator and residue
\[
 S_\alpha(c)=\partial_y(c/\eta),\qquad
 h_\alpha(f)=\frac1\eta
       \sum_{j=1}^m\frac{S_\alpha^{j-1}(c_j)}{(j-1)!}.
                                                               \tag{I.5}
\]
This explicit reduction invariant is not the ordinary coefficient of
$p^{-1}$.

\begin{lemma}[Reduction and parameter transport]
\label{int:residue}
For every finite $t$-pole $\alpha$,
\[
 h_\alpha(D_yg)=0,\qquad
 h_\alpha(D_xf)=\delta_\alpha h_\alpha(f).              \tag{I.6}
\]
All finite residues vanish if and only if
$f=D_yG+H$ for some $G\in K(t)$ and $H\in K[t]$.
\end{lemma}
\begin{proof}
Since $D_yp=\eta$, for $j\geq2$ we have
\[
 \frac{c}{p^j}
 =D_y\!\left(-\frac{c}{(j-1)\eta p^{j-1}}\right)
       +\frac{S_\alpha(c)}{(j-1)p^{j-1}}.              \tag{I.7}
\]
Repeated reduction leaves exactly $\eta h_\alpha(f)/p$.
The residue of $D_y(c/p^j)$ vanishes by cancellation of its two
successive terms in (I.5).  Distinct finite poles and the polynomial
part are independent, proving the final assertion.

It suffices to verify transport on a simple reduced term, since
$D_xD_y=D_yD_x$.  Direct calculation gives
\[
 h_\alpha\!\left(D_x\frac{c}{p}\right)
   =\frac{c_x+\partial_y(vc)}{\eta}
   =\delta_\alpha\!\left(\frac c\eta\right).          \tag{I.8}
\]
The second equality follows from
\[
 \frac{\delta_\alpha\eta}{\eta}+\partial_yv=0,
 \qquad \eta_x=-\nu_y.                                \tag{I.9}
\]
Dividing the reduced coefficient by $\eta$ is essential: without this
normalization the transport has a zeroth-order term and is not a
derivation.  Equations (I.6) imply, for every $P\in k[D_x]$,
$h_\alpha(Pf)=P(\delta_\alpha)h_\alpha(f)$.
\end{proof}

\begin{theorem}[Complete decision for one primitive extension]
\label{int:primitive}
For every rational $f\in F$ in (I.3), a telescope with coefficients in
$k$ and certificate in $F$ exists if and only if each normalized residue
$h_\alpha(f)$ is $D$-finite over $k$ for $\delta_\alpha$.
This condition is decidable by Corollary~\ref{int:bound} and the finite
tests (I.2).  In the affirmative case a telescope and its certificate
can be constructed, with a computable finite order bound.
\end{theorem}
\begin{proof}
Necessity follows immediately from Lemma~\ref{int:residue}.
For sufficiency, construct an annihilator over $k$ for each residue and
take a common left multiple $P\in k[D_x]$.
The lemma gives $Pf=D_yG+H$ with $G\in K(t)$ and $H\in K[t]$.
Average over $\operatorname{Gal}(K/k(y))$.  The original $Pf$ is fixed
and the Galois action commutes with differentiation, so the averaged
identity has $G\in F$ and $H\in k(y)[t]$.

We prove by induction on $n=\deg_tH$ that $H$ telescopes with a
certificate in $k(y)[t]$.  Write $H=h_nt^n+\cdots$.
Rational creative telescoping gives $P_0\in k[D_x]\setminus\{0\}$
and $g_n\in k(y)$ such that $P_0h_n=\partial_yg_n$.
Since $D_xt=b$ lowers the degree of a power of $t$, the polynomial
\[
             P_0H-D_y(g_nt^n)
\]
has degree at most $n-1$: its leading coefficient vanishes, and all
terms involving derivatives of $t$ have lower degree.
By induction it has a telescope $Q_0$ and polynomial
certificate $J_0$.  Thus
\[
       Q_0P_0H=D_y\bigl(Q_0(g_nt^n)+J_0\bigr).
\]
The degree-zero case is rational creative telescoping itself.
We obtain $QH=D_yJ$ and hence $QPf=D_y(QG+J)$ in $F$.

For the order bound, each transcendental residue contributes at most
$2[K:k(h_\alpha)]$ to the order of the common left multiple; algebraic
residues contribute their finite extension degrees.  Each subsequent
rational-telescoping stage has a finite effective bound, for example
the degree of the squarefree denominator in $y$, and there are finitely
many stages.  These computations give a finite bound before the final
search for a minimal telescope.  They do not give a compact bound solely
in terms of the original displayed numerator degrees.
\end{proof}

\subsection{Checks, scope, and prior work}

The symbolic verifier uses exact rational arithmetic.  For $a=1/y$ and
$\alpha^2=y+x$, set $z=\alpha$, so $y=z^2-x$.  The rational input
\[
 f=\frac{-t+2(y+x)/y}{t^2-y-x}
\]
has normalized residue $z$ at $t=z$, and its transported derivation is
\[
             \delta_\alpha z=\frac1{2z-z^2+x}.
\]
Here $K=k(z)$ and $[K:k(z)]=1$.  The three functions
$z,\delta_\alpha z,\delta_\alpha^2z$ are independent over $k$.
The sharp regular-field stopping bound $d_0+1=2$ therefore proves that
this input has no telescope in the specified field.

For a positive coupled example take $a=1+1/y$, $p=t-x-y$, and
$f=(1+y^{-3})/p$.  Its normalized residue is $y+y^{-2}$ and
$\delta=\partial_x+y\partial_y$.  Its exact minimal order is two:
\[
 (D_x-1)(D_x+2)f=D_yg,\qquad
 g=-(D_x+2)\frac y p-(D_x-1)\frac{y^{-2}}p.
\]
The verifier also checks (I.9), simple- and double-pole transport for
$\alpha=xy$ and $\alpha=x+y^2$, implicit differentiation for the
algebraic pole, and an explicit three-step polynomial-remainder
certificate.  For $t=\log(x+y)$ it verifies $D_x(1/t)=D_y(1/t)$.
For $t=\log(x^2+y^2)$, the residue of $1/t$ is
$(x^2+y^2)/(2y)$ and transport is $\partial_x-(x/y)\partial_y$;
the four derivatives required by its degree-two bound are independent.
It additionally verifies transport at a coupled pole in each of these
two parameter-dependent primitive fields.
These checks supplement the proofs; no uncontrolled
finite prefix is used to infer nonexistence.

The classical inputs must be separated from the present application.
Seidenberg's \emph{Derivations and integral closure} (1966)~\cite{Seidenberg1966} supplies the
normalization theorem.  Ngo--Nguyen--van der Put--Top~\cite{Ngo2015},
\emph{Equivalence of differential equations of order one} (2015),
Theorem~2.1 and Remark~2.2, identify regular curve derivations with the
Painlev\'e property and give an effective regularity test, building on
Muntingh--van der Put (2007).  Matzat--van der Put (2003),
Proposition~2.10, characterizes $D$-finite elements inside a
Picard--Vessiot field by membership in its Picard--Vessiot ring.
The argument here uses these established geometric ideas to obtain the
explicit field-degree bound and the normalized-residue application.

Combot's \emph{Symbolic integration on planar differential foliations}~\cite{Combot2023}
(arXiv:2306.12573, 2023), Section~4.1, already gives a terminating
algorithm for rational $G(y,\log y+h)$; it includes earlier weighted
Cauchy-kernel examples.  The present input $f(x,y,t)$ permits independent
parameter dependence and arbitrary coupled algebraic poles.  Complete
reduction in primitive towers~\cite{Du2025} and in a transcendental
Liouvillian extension~\cite{CDG2026} supplies related residue and
bounded-order frameworks.  No exhaustive priority claim is made here.
In particular, Theorem~\ref{int:primitive} is a complete result for its
stated single-extension class, and is not a complete resolution of the
original general CK6 problem.

\section{The sharp exponent window in Qi--Lim--Nantomah Open Problem 3}
\label{beta:section}

Let $\psi=\Gamma'/\Gamma$ denote the digamma function. In Open Problem~3
of \cite[Section~4]{QLN2025}, Qi, Lim, and Nantomah ask for the largest set
of real $\beta$ such that
\begin{equation}\label{beta:original}
 F_\beta(x):=x^\beta\bigl[\psi(x)+x\psi'(x)\bigr]
                 -\psi(x)\psi(1/x)>0\qquad(x>0).
\end{equation}
The result below preserves the original function, its entire positive domain,
and its strict inequality. The case $\beta=0$ was established in
\cite{QLN2025}; see also \cite{SLQ2026}. Our literature search did not locate
the sharp range proved here. This search does not certify historical priority.
In particular, related results and a different exponent problem from the same
source already appear in \cite{Pudim}; those results are not claimed here.

Put
\begin{equation}\label{beta:HP}
 H(x)=\psi(x)+x\psi'(x),\qquad P(x)=\psi(x)\psi(1/x),
\end{equation}
and denote the unique positive zero of $\psi$ by $\lambda$.
For $1/\lambda<x<1$, define
\begin{equation}\label{beta:BC}
 B(x)=\frac{\log(P(x)/H(x))}{\log x},\qquad
 C(x)=x\log x\left(\frac{P'(x)}{P(x)}-\frac{H'(x)}{H(x)}\right)
                          -\log\frac{P(x)}{H(x)}.
\end{equation}
The positivity needed for these definitions is proved below.

\begin{theorem}\label{beta:main}
The function $C$ has exactly one zero $x_*$ in $(1/\lambda,1)$.
With $\beta_*=B(x_*)$, the full answer to \eqref{beta:original} is
\begin{equation}\label{beta:answer}
 \bigl\{\beta\in\mathbb R:F_\beta(x)>0\text{ for every }x>0\bigr\}
       =[-1,\beta_*).
\end{equation}
The endpoint is specified by a unique equation and has the rigorous enclosures
\begin{align}
 \frac{76859975974774087}{10^{17}}
       &<x_*<\frac{76859975974774089}{10^{17}},\label{beta:rootbox}\\
 5.9728368638449725&<\beta_*<5.9728368638450275.
 \label{beta:thresholdbox}
\end{align}
The upper endpoint is excluded: $F_{\beta_*}(x_*)=0$.
\end{theorem}

The proof is analytic except for the compact positivity lemma
\ref{beta:certified-lemma}, whose explicit expression is certified on a finite
partition by rigorous ball arithmetic. The full verification code is included
with this paper. Thus the conclusion is unconditional, with a stated
computer-assisted component; it does not rely on numerical sampling.

\subsection{A global admissible exponent}

We use the classical bound
\begin{equation}\label{beta:productbound}
 P(x)\leq\gamma^2\qquad(x>0),
\end{equation}
where $\gamma$ is Euler's constant; see
\cite[Proposition~4]{AlzerJameson2017} and \cite[Theorem~2]{QLN2025}.
The elementary inequalities
\begin{equation}\label{beta:constants}
 0<\gamma<\frac35,\qquad \frac43<\lambda<\frac32
\end{equation}
will suffice. For the first upper bound, strict convexity of $u\mapsto1/u$
on the intervals $[k-1/2,k+1/2]$, $k\geq2$, gives
$\gamma<1-\log(3/2)<3/5$; the last inequality follows from
$\log(3/2)>2/5$. The two digamma signs giving the bounds on $\lambda$
are also checked by the certificate below.

The recurrence and positive series for the polygamma functions give
\begin{align}
 H(x)&=-\gamma+\sum_{k=1}^\infty
                 \frac{x(2k+x)}{k(k+x)^2},\label{beta:Hseries}\\
 H'(x)&=\sum_{k=1}^\infty\frac{2k}{(k+x)^3}>0.
 \label{beta:Hprime}
\end{align}
Consequently $H$ increases strictly from $-\gamma$ to infinity and has
exactly one zero, denoted by $h$. Keeping the first three positive summands
in \eqref{beta:Hseries} yields
\begin{align}
 H(1/3)&>-\gamma+\frac7{16}+\frac{13}{98}+\frac{19}{300}>0,
 \label{beta:hbound}\\
 H(2/3)&>-\gamma+\frac{16}{25}+\frac7{32}+\frac{40}{363}
             >\frac{24}{25}-\gamma>\gamma^2.
 \label{beta:Hbound}
\end{align}
In particular $0<h<1/3<1/\lambda$.
We also use
\begin{equation}\label{beta:trigamma}
 \psi'(x)>\frac1x+\frac1{2x^2}\qquad(x>0).
\end{equation}
Indeed, the trapezoidal rule strictly overestimates the integral of the
strictly convex function $u\mapsto(x+u)^{-2}$ on each $[k,k+1]$.
Summing over $k\geq0$ proves \eqref{beta:trigamma}.

\begin{lemma}\label{beta:minusone}
For every $x>0$, $F_{-1}(x)>0$.
\end{lemma}
\begin{proof}
We prove $H(x)/x>P(x)$ in five intervals, including their endpoints where
appropriate.

If $0<x<h$, then $x<1/3$ and $x+1<4/3<\lambda$. Hence
$\psi(x)=\psi(x+1)-1/x<-1/x$, while
$\psi(1/x)>\psi(3)=3/2-\gamma>\gamma$. Since $H(x)>-\gamma$,
\[
 P(x)<-\frac\gamma x<\frac{H(x)}x.
\]
If $h\leq x\leq1/\lambda$, then $H(x)\geq0$ and $P(x)\leq0$;
they never vanish simultaneously, so the required inequality is strict.

If $1/\lambda\leq x\leq1$, then $x>2/3$ and
\eqref{beta:Hbound}--\eqref{beta:productbound} imply
\[
 \frac{H(x)}x\geq H(x)>\gamma^2\geq P(x).
\]
If $1\leq x\leq\lambda<3/2$, use $\psi(x)\geq-\gamma$ and
\eqref{beta:trigamma} to obtain
\[
 \frac{H(x)}x>
 \frac{1-\gamma}{x}+\frac1{2x^2}
 >\frac{22}{45}>\frac9{25}>\gamma^2\geq P(x).
\]
Finally, if $x\geq\lambda$, then $H(x)>0$ and $P(x)\leq0$.
This completes all five intervals.
\end{proof}

\subsection{Reduction to a uniquely determined compact minimum}

The standard expansions at zero and infinity imply, as $x\to0^+$,
\begin{equation}\label{beta:asymptotics}
 H(x)=-\gamma+O(x),\qquad
 P(x)\sim-\frac{\log(1/x)}x.
\end{equation}
If $\beta<-1$, the negative term $-\gamma x^\beta$ dominates
$\log(1/x)/x$, so \eqref{beta:original} fails near zero. Therefore
$\beta\geq-1$ is necessary.

Suppose now that $\beta\geq-1$. On $(0,h)$, increasing $\beta$ from
$-1$ increases $x^\beta H(x)$, since $x<1$ and $H(x)<0$.
On $[h,1/\lambda]$, the signs $H\geq0$ and $P\leq0$ suffice.
On $[1,\infty)$, increasing $\beta$ also increases $x^\beta H(x)$,
since $H(x)>0$. Lemma~\ref{beta:minusone} therefore proves the desired
strict positivity outside $(1/\lambda,1)$.
On that remaining interval, $H>P>0$, by \eqref{beta:Hbound} and
\eqref{beta:productbound}. Taking logarithms and using $\log x<0$ gives
\begin{equation}\label{beta:compactreduction}
 F_\beta(x)>0\quad\Longleftrightarrow\quad\beta<B(x).
\end{equation}

For the uniqueness argument put
\[
 t=-\log x,\qquad T=\log\lambda,\qquad
 W(t)=\log\frac{H(e^{-t})}{P(e^{-t})}\quad(0<t<T).
\]
Then $B(e^{-t})=W(t)/t$. At the first endpoint,
\[
 W(0)=\log\frac{\zeta(2)-\gamma}{\gamma^2}>0,
\]
where the expression denotes the continuous extension to $t=0$.
At the other endpoint, the simple zero of $\psi(1/x)$ gives
\begin{equation}\label{beta:Wend}
 W(t)=-\log(T-t)+O(1),\qquad
 W'(t)=\frac1{T-t}+O(1)\quad(t\to T^-).
\end{equation}

\begin{lemma}[Certified compact convexity]\label{beta:certified-lemma}
The function $W$ satisfies $W''(t)>0$ for $0<t<T$.
\end{lemma}

\begin{proof}
We specify the exact expression and its finite verification.
Write
\[
 z=\psi(x),\ u=\psi'(x),\ v=\psi''(x),\ w=\psi'''(x),\qquad
 Z=\psi(1/x),\ U=\psi'(1/x),\ V=\psi''(1/x).
\]
Thus $H=z+xu$, $H'=2u+xv$, and $H''=3v+xw$.
For $D=x\,d/dx$, direct differentiation yields
\begin{equation}\label{beta:clearedidentity}
 D^2\log\frac HP=\frac{\mathcal N(x)}{x^2H^2z^2Z^2},
\end{equation}
where
\begin{align}
 \mathcal N(x)={}&
 x^2\bigl[(xH'+x^2H'')H-x^2(H')^2\bigr]z^2Z^2\notag\\
 &-x^2\bigl[(xu+x^2v)z-x^2u^2\bigr]H^2Z^2\notag\\
 &-\bigl[(xU+V)Z-U^2\bigr]H^2z^2.
 \label{beta:numerator}
\end{align}
Unlike the logarithmic expression, $\mathcal N$ has no singularity at
$Z=0$. It can therefore be verified on the larger rational interval
$[2/3,1]$, which contains $[1/\lambda,1]$.

The accompanying script \texttt{beta\_problem3\_certify.py} evaluates
\eqref{beta:numerator} using Arb's rigorous ball arithmetic
\cite{Johansson2017}, through python-flint~0.9.0 at 192-bit precision.
Starting from $[2/3,1]$, it bisects intervals at their exact rational midpoints.
An interval is accepted only if the ball enclosing the entire expression
$\mathcal N$ is strictly positive. Digamma values are evaluated by Arb's
digamma routine, and for $n\geq1$ the identity
\[
 \psi^{(n)}(x)=(-1)^{n+1}n!\,\zeta(n+1,x)
\]
supplies the other values using rigorous Hurwitz-zeta evaluation.
The run terminates with 64 accepted intervals; all endpoints and resulting
balls are recorded in \texttt{beta\_problem3\_boxes.json}. Its relevant
output is
\begin{verbatim}
CERTIFIED cleared log-convexity positive on [2/3,1]; pieces = 64
CERTIFIED unique digamma zero in (4/3,3/2)
\end{verbatim}
This verifies $\mathcal N(x)>0$ throughout $[2/3,1]$.
On $(1/\lambda,1)$ the denominator in \eqref{beta:clearedidentity} is
strictly positive. Finally, $D=-d/dt$, so $D^2\log(H/P)=W''(t)$.
\end{proof}

\begin{proof}[Proof of Theorem~\ref{beta:main}]
By Lemma~\ref{beta:certified-lemma},
\[
 \frac d{dt}\bigl[tW'(t)-W(t)\bigr]=tW''(t)>0.
\]
The expression in brackets tends to $-W(0)<0$ as $t\to0^+$ and to
infinity as $t\to T^-$, by \eqref{beta:Wend}. Thus it has exactly one
zero. Consequently $W(t)/t$, and equivalently $B(x)$, has exactly one
minimum. Differentiating \eqref{beta:BC} gives
\[
 B'(x)=\frac{C(x)}{x(\log x)^2},
\]
which proves the asserted unique zero of $C$.

For the two exact rational endpoints $x_-,x_+$ in
\eqref{beta:rootbox}, the same script certifies
$C(x_-)<0<C(x_+)$. Evaluation of $B$ on the full ball
$[x_-,x_+]$ gives
\begin{verbatim}
CERTIFIED critical point in
76859975974774087/100000000000000000
76859975974774089/100000000000000000
CERTIFIED threshold enclosure = [5.9728368638450 +/- 2.75e-14]
\end{verbatim}
This implies \eqref{beta:rootbox}--\eqref{beta:thresholdbox}.
The necessity of $\beta\geq-1$ and the compact equivalence
\eqref{beta:compactreduction} now show that precisely
$-1\leq\beta<\min B=\beta_*$ are admissible. At the minimum,
$x_*^{\beta_*}H(x_*)=P(x_*)$, so the upper endpoint is excluded.
\end{proof}

\begin{remark}\label{beta:scope}
The theorem determines the original global parameter set, including both
endpoint decisions. The implicit constant is specified by a proved unique
equation with a certified rational isolating interval, rather than by an
unevaluated global optimization. The reduction to compact convexity and the
use of interval arithmetic are established mathematical techniques; no claim
of a historically new general proof method is made.
\end{remark}


\section{Limits and excluded claims}
The results above give finite criteria and one sharp parameter answer.
They do not establish the requested target of one hundred newly solved
difficult conjectures. The full elementary-function classification in
Chen--Kauers Problem~8 and the arbitrary-tower existence problem in
Problem~6 remain outside the proved scope.

An independently derived proof of complete monotonicity of
$[\log\Gamma(x)+\log\Gamma(1/x)]''$ was removed from the new-results
list after a valid public proof dated 17 May 2026 was located in the
Pudim repository~\cite{Pudim}. The same repository also discusses other
questions from Qi--Lim--Nantomah, including its separate $n=2$ problem.
This does not resolve Open Problem~3, whose $n=0$ parameter range is
the subject of the present theorem. The 2026 alternative-positivity
paper~\cite{SLQ2026} proves the unweighted case, not the sharp full range.

Graph and Erd\H{o}s-problem investigations in the same round produced
partial reductions and finite searches, not additional complete
closures. None is included in the theorem count. A failed bounded
counterexample search is not a proof of a universal statement.

\appendix
\section{Reproduction}
Compile this source twice with \texttt{pdflatex}. In the compilation
directory, run the extracted Python programs listed below. The compact
polygamma certificate uses \texttt{python-flint 0.9.0} at 192-bit
precision; the symbolic checks use \texttt{SymPy 1.14.0}. All accepted
intervals are certified in full, rather than sampled at points.
\subsection*{algebraic\_monodromy\_verify.py}
\lstinputlisting{algebraic_monodromy_verify.py}
\subsection*{integration\_verify.py}
\lstinputlisting{integration_verify.py}
\subsection*{beta\_problem3\_certify.py}
\lstinputlisting{beta_problem3_certify.py}
\begin{thebibliography}{99}
\bibitem{CK2016} S. Chen and M. Kauers,
\emph{Some Open Problems related to Creative Telescoping},
arXiv:1609.03768 (2016). \url{https://arxiv.org/abs/1609.03768}.

\bibitem{Chen2022} S. Chen,
\emph{Stability Problems in Symbolic Integration},
Proceedings of ISSAC 2022, 517--524.
\url{https://doi.org/10.1145/3476446.3535502};
arXiv:2202.06305, Section 5.

\bibitem{GP2014} L. Gavrilov and F. Pakovich,
\emph{Moments on Riemann surfaces and hyperelliptic Abelian integrals},
Comment. Math. Helv. 89 (2014), 125--155.
\url{https://doi.org/10.4171/CMH/314}; arXiv:1107.3029.

\bibitem{SSS2025} C. Sahu, M. Seiss, and V. R. Srinivasan,
\emph{Iterated and Generalized Iterated Integrals},
arXiv:2504.04846 (2025). \url{https://arxiv.org/abs/2504.04846}.

\bibitem{Rosenlicht1972} M. Rosenlicht,
\emph{Integration in finite terms}, Amer. Math. Monthly 79 (1972), 963--972.
\url{https://doi.org/10.1080/00029890.1972.11993166}.

\bibitem{Seidenberg1966} A. Seidenberg,
\emph{Derivations and integral closure}, Pacific J. Math. 16 (1966), 167--173.
\url{https://msp.org/pjm/1966/16-1/pjm-v16-n1-p16-s.pdf}.

\bibitem{Ngo2015} L. X. C. Ngo, K. A. Nguyen, M. van der Put, and J. Top,
\emph{Equivalence of differential equations of order one},
J. Symbolic Comput. 71 (2015), 47--59.
\url{https://doi.org/10.1016/j.jsc.2014.09.041}.

\bibitem{Combot2023} T. Combot,
\emph{Symbolic integration on planar differential foliations},
arXiv:2306.12573 (2023), Section 4.1.
\url{https://arxiv.org/abs/2306.12573}.

\bibitem{Du2025} H. Du, Y. Gao, W. Li, and Z. Li,
\emph{Complete Reduction for Derivatives in a Primitive Tower},
arXiv:2510.13456 (2025).
\url{https://arxiv.org/abs/2510.13456}.

\bibitem{CDG2026} S. Chen et al.,
\emph{Complete Reduction for Derivatives in a Transcendental Liouvillian Extension},
arXiv:2602.03592v3 (2026).
\url{https://arxiv.org/abs/2602.03592}.

\bibitem{QLN2025} F. Qi, D. Lim, and K. Nantomah,
\emph{Monotonicity and positivity of several functions involving ratios and
products of polygamma functions}, J. Inequal. Appl. 2025, article 5.
\url{https://doi.org/10.1186/s13660-024-03245-8}, Section 4, Open Problem 3.

\bibitem{AlzerJameson2017} H. Alzer and G. Jameson,
\emph{A harmonic mean inequality for the digamma function and related results},
Rend. Semin. Mat. Univ. Padova 137 (2017), 203--209.
\url{https://doi.org/10.4171/RSMUP/137-10}.

\bibitem{SLQ2026} Y. Shuang, D. Lim, and F. Qi,
\emph{Alternative proofs for the positivity of a function involving the
product of two digamma functions}, Aust. J. Math. Anal. Appl. 23 (2026),
no. 1, article 5, 6 pp.
\url{https://ajmaa.org/searchroot/files/pdf/v23n1/v23i1p5.pdf}.

\bibitem{Johansson2017} F. Johansson,
\emph{Arb: efficient arbitrary-precision midpoint-radius interval arithmetic},
IEEE Trans. Comput. 66 (2017), 1281--1292.
\url{https://doi.org/10.1109/TC.2017.2690633}.

\bibitem{Pudim} Pudim project,
\emph{pudim-ai-demo-zetalaw}, public research repository,
\texttt{APPLICATIONS.md}, APP-0004 and associated staged proof
(dated 17 May 2026); APP-0007, APP-0008, and APP-0011.
Inspected 2 November 2026.
\url{https://github.com/pudim-project/pudim-ai-demo-zetalaw}.

\bibitem{VMRC} QEEDDD,
\emph{Verifiable Mathematical Reasoning Corpus},
commit \texttt{ecd87de3348dcb706a76bb12fde722d7a348a012},
2 November 2026.
\url{https://github.com/QEEDDD/Verifiable-Mathematical-Reasoning-Corpus}.
\end{thebibliography}

\end{document}
